\begin{lstlisting}[language=Python]
# Implementação do gráfico da variância
plt.rc('lines', linewidth=2.5)
fig, ax = plt.subplots(figsize=(10, 6))

ax.set_title('Variância da Média Amostral (Lei dos Grandes Números)')
ax.set_xlabel('Número de Iterações (n)')
ax.set_ylabel('Variância da Média (sigma^2/n)')
ax.grid(True)
ax.set_xscale('log') # Escala log no eixo X para visualizar todas as iterações
ax.set_yscale('log') # Escala log no eixo Y pois a variância decai exponencialmente

# Começa no índice 10 para evitar instabilidade numérica nas primeiras iterações
start = 10

# Linha tracejada: variância teórica da média, o comportamento ideal esperado
ax.plot(nArray[start:], theoreticalVarianceOfMean[start:], dashes=[6,2], label='sigma^2/n (teórico)')

# Linha contínua: variância empírica da média calculada a cada iteração
ax.plot(nArray[start:], varianceOfMean[start:], label='variância empírica da média', alpha=0.7)

ax.legend(handlelength=4)
plt.tight_layout()
plt.show()
\end{lstlisting}
